\chapter{Routing Is Geometry}

The old enterprise stack routes through abstraction. It models supply chains and operational logistics by abstracting the world into graphs, running algorithms over those symbolic nodes, and interpreting the consequences post-facto on a dashboard. The \textbf{uni*} ecosystem upends this paradigm. We give operations a finite, physical geometry, and we route through lawful space.

\section{The Failure of Abstract Routing}

When routing is treated as an abstract graph problem floating above the physical system, the result is semantic disconnect. The process intent (e.g., ``Reroute vessel to a port with available cold-chain capacity'') must be interpreted against a fragile, downstream reconstruction of reality. 

\section{The Operational Globe}

In this paradigm, operations are mapped to a canonical globe coordinate. A route is not merely represented by geometry; a route \textbf{is} geometry. A route is defined as a trajectory through \texttt{GlobeCell(domain, cell, place)}. This unifies the geographical position with the operational state. 

\section{Route = Lawful Trajectory}

\textbf{At $64^2$, the world becomes an attention surface. At $64^3$, every point on that surface gains operational depth. Routing stops being graph abstraction and becomes lawful geometry.}

Global operations become computable when routing is treated as this lawful geometry. The subsequent architecture presented in this whitepaper is entirely in service of executing these trajectories securely, deterministically, and at nanosecond scales.